\section{Related structures}

\paragraph{Rauzy--Veech and AGY dynamics.}
Rauzy induction records interval-exchange dynamics by a directed symbolic
system whose integral matrices multiply in chronological order.  The
accelerated model of \citet{agy2006} supplies a countable Markov return map,
a logarithmic roof, and the exponential-tail estimates needed for transfer
operators.  Our source is a frozen four-letter, genus-two instance of that
framework.  We do not alter its clock, average its transition matrices, or
identify paths merely because they share a characteristic polynomial.  The
new question is what arithmetic information survives when the same integral
symplectic cocycle is reduced modulo primes.

\paragraph{Holomorphic transfer operators.}
Compactly contained holomorphic map-weight systems admit strong eigenvalue
and trace-ideal bounds on Bergman-type spaces
\citep{bandtlow2008explicit,bandtlow2008ruelle}.  The scalar companion
construction provides exactly such a common domain for the countable AGY
branches and proves locally uniform trace-norm summability.  The present
paper does not rederive that complex-cone theorem.  Instead, it asks whether
an arithmetic vector fibre preserves the scalar result.  Finite dimension is
decisive: tensoring a branch with a \(p^2\)-dimensional unitary multiplies its
trace norm by \(p^2\), whereas the unsmoothed real oscillator has infinite
multiplicity.

\paragraph{Weil representations over finite and real fields.}
The oscillator or metaplectic representation over \(\R\) is naturally tied
to symplectic phase space \citep{folland1989}.  Over a finite field of odd
characteristic, the canonical intertwining construction yields an honest
linear Weil representation of the symplectic group
\citep{gurevich2009}.  The genus-two fibre used here is the full
\(p^2\)-dimensional representation.  In particular, we do not add edgewise
metaplectic signs, quotient by the central element \(-I\), or silently retain
only an even or odd parity constituent.

\paragraph{Exact characters and dynamical weights.}
\citet{thomas2008} expresses the finite Weil character through the fixed
space of \(g\) and the discriminant of a bilinear form on
\(V/\ker(g-I)\).  That theorem makes singular primes computable without a
dense \(p^2\)-by-\(p^2\) matrix.  Our use differs from a pointwise character
Euler product: the character is evaluated on a genuine chronological product,
and repetitions use \(\Theta_p(g^r)\), not \(\Theta_p(g)^r\).  We also
reconstruct the complete characteristic polynomial of the fibre operator by
Newton identities.

\paragraph{Relation to Hilbert--P\'olya searches.}
A trace-class dynamical determinant and a natural unitary fibre address only
two structural requirements of a Hilbert--P\'olya programme.  A successful
global construction would still need an intrinsic prime organization, a
controlled analytic continuation, the appropriate functional equation and
gamma factor, and a bridge to a self-adjoint spectral problem.  The present
work treats the absence of those features as a result, not as a normalization
detail.  Its main negative controls quantify information lost by every
class-function fibre before any comparison with Riemann zeros is attempted.
